\documentclass[11pt]{article}
\usepackage[margin=1in]{geometry}
\usepackage{amsmath,amssymb}
\usepackage{graphicx}
\usepackage[hidelinks]{hyperref}

\title{Literature Status: Positivity of the Real \texorpdfstring{$\ell_p$}{ell-p} Numerical Index}
\author{FA/Banach run \texttt{fa\_banach\_001}}
\date{2026-06-14}

\begin{document}
\maketitle

\section*{Status}

This is a literature-already-answered packet, not an original proof from the
run.  Problem 1 in Kadets--Martin--Paya, arXiv:math/0605781, asks whether
the numerical index $n(\ell_p)$ is positive for every real $\ell_p$ with
$p\neq 2$.  Martin--Meri--Popov, arXiv:0903.2704, answer this affirmatively
in the stronger form
\[
  n(L_p(\mu))>0
  \qquad (1<p<\infty,\ p\neq 2)
\]
for every positive measure $\mu$ in the real case.

\section*{Original Problem}

The original source is:

\begin{quote}
V. Kadets, M. Martin, R. Paya, \emph{Recent progress and open questions on
the numerical index of Banach spaces}, arXiv:math/0605781.
\end{quote}

After recalling that
\[
 n(L_p[0,1])=n(\ell_p)=\inf_m n(\ell_p^m),
\]
the paper states that $n(\ell_p^m)>0$ for each finite $m$ and each
$p\neq 2$, but that the infinite-dimensional positivity was unknown:

\begin{center}
\includegraphics[width=\textwidth]{figures/open_problem_crop.png}
\end{center}

\section*{Later Answer}

The supporting source is:

\begin{quote}
M. Martin, J. Meri, M. Popov, \emph{On the numerical index of real
$L_p(\mu)$-spaces}, arXiv:0903.2704.
\end{quote}

Corollary 3 gives the quantitative lower bound
\[
  n(L_p(\mu))\geq \frac{M_p}{12e},
  \qquad
  M_p=\max_t \frac{|t^{p-1}-t|}{1+t^p}.
\]
Since $M_p>0$ when $p\neq 2$, Corollary 4 states the positive answer for
all real $L_p(\mu)$ spaces, $1<p<\infty$, $p\neq 2$, and explicitly says
that this answers the question also posed in reference [7, Problem 1].
Reference [7] is the 2006 Kadets--Martin--Paya survey.

\begin{center}
\includegraphics[width=\textwidth]{figures/supporting_answer_crop.png}
\end{center}

\section*{Proof Intuition}

The supporting result does more than the original $\ell_p$ question asks:
it proves a uniform numerical-radius lower bound for every operator on every
real $L_p(\mu)$.  The constant is positive exactly outside the Hilbertian
case $p=2$.  Taking $\mu$ to be counting measure on $\mathbb N$ specializes
the theorem to $\ell_p$, so the original problem follows immediately.

\section*{Verification Notes}

The answer is not contained in the original 2006 paper; there the positivity
of $n(\ell_p)$ is posed as Problem 1.  The supporting paper is a separate
later paper and explicitly cites the 2006 survey as [7].  Its Corollary 4
answers [7, Problem 1] in the positive.

The endpoint cases $p=1$ and $p=\infty$ were already standard in the 2006
source, because $L$-spaces and $M$-spaces have numerical index $1$.  Thus
the supporting theorem resolves the only nontrivial range of the question.

\section*{Search Bounds}

The local registry, solution index, and attempt index were checked for
arXiv:0605781, arXiv:0903.2704, positivity of $n(\ell_p)$, and the adjacent
two-dimensional formula.  A web search for numerical-index positivity of
$\ell_p$ and $L_p$ spaces located arXiv:0903.2704.  The arXiv page, source
paper, and supporting PDF were inspected.

\section*{Limitations}

This packet does not claim an original proof.  It also does not address the
neighboring 2006 questions asking whether $n(\ell_p)=n(\ell_p^2)$ or asking
for the complex two-dimensional formula.

\section*{Human Review Recommendation}

Check two points: first, that [7] in arXiv:0903.2704 is indeed the
Kadets--Martin--Paya survey; second, that the specialization of real
$L_p(\mu)$ to $\ell_p$ via counting measure matches the convention in the
original problem.

\end{document}
